
    \documentclass{standalone}
    \usepackage[english,french]{babel} 
    \usepackage{tikz}
    \usetikzlibrary{calc}
    \usepackage[T1]{fontenc}
    \usepackage[utf8]{inputenc}         
    \usepackage{pgfplotstable}
    \usepackage{pgfplots}
    \pgfplotsset{compat=1.18}
    \usepackage{siunitx}               
    \usepackage{booktabs} 
    \begin{document}
    \begin{tikzpicture}
    \begin{axis}[
                height=10cm,
                width=15cm,
                tick align=outside,
                tick pos=left,
                ylabel={Proportion de la catégorie avec données manquante[\%]},
                ymin= 0,
                ymax= 100,
                ytick={0,20,40,60,80,100},
                yticklabels={0,20,40,60,80,100},
                legend style={
                    at={(0.5,1.02)},  % x=0..1, y=0..1 relative to axis
                    anchor=south      % which part of the legend is placed at that point
                },
                x tick label style={rotate=45, anchor=east},
                xtick={1,2,3,4,5,6,7,8,9,10},
                xticklabels={Usage mixte,
Résidentielle,
Industrielle,
Infrastructure,
Commerciale,
Services,
Récréation,
Ressources naturelles,
Inexploité,
Sans Rôle},
                xlabel={Utilisation du sol},
                xmin=-0.5,
                xmax=10.5
            ]\addplot[ybar,draw=none,fill=blue,legend image post style={/pgfplots/ybar legend},bar width= 8pt,bar shift=-4pt]
                coordinates{
                    (1,36.58536585365854)(2,0.37036153962312135)(3,48.5632183908046)(4,86.8037703513282)(5,34.56221198156682)(6,30.593220338983052)(7,75.68)(8,79.38931297709924)(9,97.06469079509831)(10,nan)
                };
            \addlegendentry{En nombre de lots};\addplot[ybar,draw=none,fill=orange,legend image post style={/pgfplots/ybar legend},bar width= 8pt,bar shift=4pt]
                coordinates{
                    (1,77.29676210547714)(2,1.4190601297424188)(3,65.01352732339161)(4,87.344971361682)(5,53.09431370368725)(6,67.30013730462127)(7,77.06616473206749)(8,82.76065603782936)(9,93.06007452581294)(10,nan)
                };
            \addlegendentry{En valeur foncière};
        \end{axis}
        \end{tikzpicture}
        \end{document} 